\documentclass[12pt,a4paper]{article}

% ─── Paquetes ───────────────────────────────────────────────────────────────
\usepackage[utf8]{inputenc}
\usepackage[T1]{fontenc}
\usepackage[spanish]{babel}
\usepackage{amsmath,amssymb}
\usepackage{booktabs}
\usepackage{array}
\usepackage{geometry}
\usepackage{xcolor}
\usepackage{titlesec}
\usepackage{enumitem}
\usepackage{fancyhdr}
\usepackage{mdframed}

% ─── Geometría ──────────────────────────────────────────────────────────────
\geometry{left=2.5cm, right=2.5cm, top=3cm, bottom=3cm}

% ─── Colores ────────────────────────────────────────────────────────────────
\definecolor{primario}{RGB}{0,90,160}
\definecolor{secundario}{RGB}{220,235,250}
\definecolor{acento}{RGB}{200,50,50}

% ─── Encabezado / pie de página ─────────────────────────────────────────────
\pagestyle{fancy}
\fancyhf{}
\fancyhead[L]{\small\textcolor{primario}{\textbf{Lógica Proposicional}}}
\fancyhead[R]{\small\textcolor{primario}{Cuarto Grado de Secundaria}}
\fancyfoot[C]{\thepage}
\renewcommand{\headrulewidth}{0.6pt}

% ─── Estilos de título ───────────────────────────────────────────────────────
\titleformat{\section}{\large\bfseries\color{primario}}{}{0pt}{}[\titlerule]
\titleformat{\subsection}{\normalsize\bfseries\color{acento}}{}{0pt}{}

% ─── Entorno para definiciones ───────────────────────────────────────────────
\newmdenv[
  backgroundcolor=secundario,
  linecolor=primario,
  linewidth=1.5pt,
  leftmargin=0pt,
  rightmargin=0pt,
  innerleftmargin=10pt,
  innerrightmargin=10pt,
  innertopmargin=8pt,
  innerbottommargin=8pt
]{definicion}

% ─── Comando para tabla de verdad ────────────────────────────────────────────
\newcommand{\TV}[1]{\textbf{#1}}   % V / F en negrita

% ─── Documento ───────────────────────────────────────────────────────────────
\begin{document}

% Portada
\begin{center}
  {\LARGE\bfseries\color{primario} CONECTIVOS LÓGICOS Y TABLA DE VERDAD}\\[4pt]
  {\large Teoría — Cuarto Grado de Secundaria}\\[2pt]
  \textcolor{gray}{\rule{0.6\linewidth}{0.6pt}}
\end{center}

\vspace{6pt}

% ════════════════════════════════════════════════════════════════════════════
\section{Lógica Proposicional}
% ════════════════════════════════════════════════════════════════════════════

\begin{definicion}
La \textbf{lógica proposicional} es una rama de la lógica que tiene por objeto de
estudio las \textbf{proposiciones} y la relación entre ellas, así como la función que
desempeñan las variables proposicionales y los conectivos lógicos.
\end{definicion}

% ════════════════════════════════════════════════════════════════════════════
\section{Proposición Lógica}
% ════════════════════════════════════════════════════════════════════════════

\begin{definicion}
Una \textbf{proposición lógica} (o enunciado cerrado) es un enunciado que posee un
valor de verdad definido: \textbf{Verdadero (V)} o \textbf{Falso (F)}.  Se denomina
\textbf{variable proposicional} y se representa con letras minúsculas ($p$, $q$, $r$, \ldots).
\end{definicion}

\subsection{Clasificación}

\begin{itemize}[leftmargin=1.5cm]
  \item \textbf{Proposiciones simples (atómicas):} No contienen conectivos lógicos.
        \begin{itemize}
          \item Carlos es despistado.
          \item Carlos es travieso.
        \end{itemize}
  \item \textbf{Proposiciones compuestas (moleculares):} Se forman uniendo proposiciones
        simples mediante conectivos lógicos.
        \begin{itemize}
          \item Carlos es travieso \textit{y} despistado.
          \item \textit{Es falso que} Daniel sea actor de cine.
        \end{itemize}
\end{itemize}

% ════════════════════════════════════════════════════════════════════════════
\section{Conectivos Lógicos}
% ════════════════════════════════════════════════════════════════════════════

A continuación se presentan los seis conectivos lógicos fundamentales con su
tabla de verdad correspondiente.

% ─── a) Conjunción ──────────────────────────────────────────────────────────
\subsection{a) Conjunción}

\begin{minipage}{0.55\linewidth}
\begin{itemize}[leftmargin=*]
  \item \textbf{Conectiva:} y / pero / e / sin embargo \ldots
  \item \textbf{Operador:} $\wedge$ (punto $\cdot$)
  \item \textbf{Forma:} $p \wedge q$
  \item \textbf{Regla:} es \textbf{V} solo cuando \emph{ambas} proposiciones son V.
\end{itemize}
\end{minipage}%
\hfill
\begin{minipage}{0.38\linewidth}
\centering
\begin{tabular}{cc|c}
\toprule
$p$ & $q$ & $p \wedge q$ \\
\midrule
\TV{V} & \TV{V} & \TV{V} \\
\TV{V} & F      & F      \\
F      & \TV{V} & F      \\
F      & F      & F      \\
\bottomrule
\end{tabular}
\end{minipage}

\bigskip

% ─── b) Disyunción débil ─────────────────────────────────────────────────────
\subsection{b) Disyunción Débil (Inclusiva)}

\begin{minipage}{0.55\linewidth}
\begin{itemize}[leftmargin=*]
  \item \textbf{Conectiva:} o / u / \ldots o \ldots
  \item \textbf{Operador:} $\vee$
  \item \textbf{Forma:} $p \vee q$
  \item \textbf{Regla:} es \textbf{F} solo cuando \emph{ambas} proposiciones son F.
\end{itemize}
\end{minipage}%
\hfill
\begin{minipage}{0.38\linewidth}
\centering
\begin{tabular}{cc|c}
\toprule
$p$ & $q$ & $p \vee q$ \\
\midrule
\TV{V} & \TV{V} & \TV{V} \\
\TV{V} & F      & \TV{V} \\
F      & \TV{V} & \TV{V} \\
F      & F      & F      \\
\bottomrule
\end{tabular}
\end{minipage}

\bigskip

% ─── c) Disyunción fuerte ────────────────────────────────────────────────────
\subsection{c) Disyunción Fuerte (Exclusiva)}

\begin{minipage}{0.55\linewidth}
\begin{itemize}[leftmargin=*]
  \item \textbf{Conectiva:} o \ldots o / o bien \ldots o bien \ldots
  \item \textbf{Operador:} $\Delta$ (o $\not\leftrightarrow$)
  \item \textbf{Forma:} $p \;\Delta\; q$
  \item \textbf{Regla:} es \textbf{V} cuando las proposiciones tienen \emph{distinto} valor.
\end{itemize}
\end{minipage}%
\hfill
\begin{minipage}{0.38\linewidth}
\centering
\begin{tabular}{cc|c}
\toprule
$p$ & $q$ & $p \;\Delta\; q$ \\
\midrule
\TV{V} & \TV{V} & F      \\
\TV{V} & F      & \TV{V} \\
F      & \TV{V} & \TV{V} \\
F      & F      & F      \\
\bottomrule
\end{tabular}
\end{minipage}

\bigskip

% ─── d) Condicional ──────────────────────────────────────────────────────────
\subsection{d) Condicional (Implicación)}

\begin{minipage}{0.55\linewidth}
\begin{itemize}[leftmargin=*]
  \item \textbf{Conectiva:} Si \ldots entonces / por lo tanto
  \item \textbf{Operador:} $\rightarrow$ (o $\Rightarrow$)
  \item \textbf{Forma:} $p \rightarrow q$
  \item \textbf{Regla:} es \textbf{F} solo cuando el antecedente es V y el consecuente es F.
\end{itemize}
\end{minipage}%
\hfill
\begin{minipage}{0.38\linewidth}
\centering
\begin{tabular}{cc|c}
\toprule
$p$ & $q$ & $p \rightarrow q$ \\
\midrule
\TV{V} & \TV{V} & \TV{V} \\
\TV{V} & F      & F      \\
F      & \TV{V} & \TV{V} \\
F      & F      & \TV{V} \\
\bottomrule
\end{tabular}
\end{minipage}

\bigskip

% ─── e) Bicondicional ────────────────────────────────────────────────────────
\subsection{e) Bicondicional (Doble Implicación)}

\begin{minipage}{0.55\linewidth}
\begin{itemize}[leftmargin=*]
  \item \textbf{Conectiva:} Si y solo si / entonces y solo entonces
  \item \textbf{Operador:} $\leftrightarrow$ (o $\equiv$)
  \item \textbf{Forma:} $p \leftrightarrow q$
  \item \textbf{Regla:} es \textbf{V} cuando ambas proposiciones tienen el \emph{mismo} valor.
\end{itemize}
\end{minipage}%
\hfill
\begin{minipage}{0.38\linewidth}
\centering
\begin{tabular}{cc|c}
\toprule
$p$ & $q$ & $p \leftrightarrow q$ \\
\midrule
\TV{V} & \TV{V} & \TV{V} \\
\TV{V} & F      & F      \\
F      & \TV{V} & F      \\
F      & F      & \TV{V} \\
\bottomrule
\end{tabular}
\end{minipage}

\bigskip

% ─── f) Negación ─────────────────────────────────────────────────────────────
\subsection{f) Negación}

\begin{minipage}{0.55\linewidth}
\begin{itemize}[leftmargin=*]
  \item \textbf{Conectiva:} no / ni / no es el caso que
  \item \textbf{Operador:} $\sim$ (o $\neg$)
  \item \textbf{Forma:} $\sim p$
  \item \textbf{Regla:} invierte el valor de verdad de la proposición.
\end{itemize}
\end{minipage}%
\hfill
\begin{minipage}{0.38\linewidth}
\centering
\begin{tabular}{c|c}
\toprule
$p$ & $\sim p$ \\
\midrule
\TV{V} & F      \\
F      & \TV{V} \\
\bottomrule
\end{tabular}
\end{minipage}

% ════════════════════════════════════════════════════════════════════════════
\section{Tablas de Valores de Verdad (Matriz Principal)}
% ════════════════════════════════════════════════════════════════════════════

\begin{definicion}
\textbf{Evaluar} un esquema molecular es obtener su \textbf{matriz principal}, es decir,
la columna de valores de verdad del conectivo principal.
\end{definicion}

\subsection{Reglas importantes}

\begin{enumerate}[leftmargin=1.5cm]
  \item El número de filas en la tabla es $2^{n}$, donde $n$ es el número de variables
        proposicionales distintas.
  \item Se debe \textbf{jerarquizar} los conectivos antes de evaluar:
        \begin{enumerate}[label=\arabic*$^\circ$]
          \item Negación ($\sim$)
          \item Conjunción ($\wedge$) y disyunción fuerte ($\Delta$)
          \item Disyunción débil ($\vee$)
          \item Condicional ($\rightarrow$)
          \item Bicondicional ($\leftrightarrow$)
          \item Los paréntesis \textbf{rompen} la jerarquía.
        \end{enumerate}
\end{enumerate}

\subsection{Ejemplo resuelto}

Obtener la matriz principal de $(p \wedge q) \rightarrow (p \;\Delta\; q)$:

\[
\begin{array}{cc|c|c|c}
\toprule
p & q & (p \wedge q) & (p \;\Delta\; q) & (p \wedge q) \rightarrow (p \;\Delta\; q) \\
  &   & \text{(paso 1)} & \text{(paso 2)} & \textbf{(paso 3 = matriz principal)} \\
\midrule
V & V & V & F & \mathbf{F} \\
V & F & F & V & \mathbf{V} \\
F & V & F & V & \mathbf{V} \\
F & F & F & F & \mathbf{V} \\
\bottomrule
\end{array}
\]

\textbf{Matriz principal:} F, V, V, V

% ════════════════════════════════════════════════════════════════════════════
\section{Clases de Matrices Principales}
% ════════════════════════════════════════════════════════════════════════════

\begin{center}
\begin{tabular}{>{\bfseries}lll}
\toprule
Tipo & Característica & Ejemplo \\
\midrule
Tautología    & Todos los valores son \textbf{V} & $p \vee \sim p$ \\
Contradicción & Todos los valores son \textbf{F} & $p \wedge \sim p$ \\
Contingente   & Hay al menos un V y un F         & $p \rightarrow q$ \\
\bottomrule
\end{tabular}
\end{center}

\vspace{4pt}

\begin{definicion}
\textbf{Resumen:}\\[4pt]
\begin{tabular}{@{}ll@{}}
$p \wedge q$           & V solo si ambas son V.\\
$p \vee q$             & F solo si ambas son F.\\
$p \;\Delta\; q$       & V si tienen \emph{distinto} valor.\\
$p \rightarrow q$      & F solo si $p$=V y $q$=F.\\
$p \leftrightarrow q$  & V si tienen el \emph{mismo} valor.\\
$\sim p$               & Invierte el valor de $p$.\\
\end{tabular}
\end{definicion}

\end{document}
